\sezione{6}{Importanza delle feature}{Parte I · MDI e importanza per permutazione}

\begin{inbreve}
Un Random Forest non ha coefficienti come una regressione, ma può dire \hl{quanto conta ogni variabile}. Ci sono due misure: la \key{MDI} (diminuzione media di impurità, \texttt{feature\_importances\_}), veloce ma distorta, e l'\key{importanza per permutazione}, più affidabile, calcolata su dati non usati per l'addestramento.
\end{inbreve}

\subsection{MDI: Mean Decrease in Impurity}

\begin{tecnica}[Definizione]
Ogni volta che un nodo $t$ viene diviso sulla feature $j$, l'impurità cala di
\[
\Delta i(t)=\frac{n_t}{n}\Big[H(t)-\frac{n_{t}^{\mathrm{sx}}}{n_t}H(t^{\mathrm{sx}})-\frac{n_t^{\mathrm{dx}}}{n_t}H(t^{\mathrm{dx}})\Big].
\]
L'importanza MDI della feature $j$ somma questi cali su tutti i nodi che usano $j$ e ne fa la media sugli alberi:
\[
\mathrm{MDI}(j)=\frac1B\sum_{b=1}^{B}\ \sum_{t\in T_b:\ v(t)=j}\Delta i(t),
\]
poi le importanze vengono normalizzate in modo da sommare a 1. In scikit-learn è l'attributo \texttt{feature\_importances\_}.
\end{tecnica}

\begin{attenzione}[I due difetti della MDI]
La documentazione di scikit-learn segnala che la MDI:
\begin{itemize}
  \item è calcolata sui \hl{dati di addestramento}, quindi può dare importanza a feature che servono solo a memorizzare il rumore;
  \item è \hl{distorta verso le feature con molti valori distinti} (alta cardinalità, ad esempio un identificativo numerico), che offrono più soglie candidate.
\end{itemize}
Per conclusioni affidabili usa l'importanza per permutazione.
\end{attenzione}

\subsection{Importanza per permutazione}

\begin{intuizione}[Mescola una colonna e guarda cosa succede]
Prendi il modello già addestrato e un set di dati di validazione. Mescola a caso \hl{solo la colonna $j$}, lasciando invariate le altre: il legame tra quella variabile e il target è distrutto. Se il punteggio del modello crolla, la variabile era importante; se non cambia, il modello non la usava (o poteva farne a meno).
\end{intuizione}

\begin{tecnica}[Definizione (sklearn.inspection)]
Sia $s$ il punteggio del modello sui dati di riferimento $D$. Per ogni feature $j$ e per ogni ripetizione $k=1,\dots,K$ si permuta la colonna $j$ ottenendo $\tilde D_{k,j}$ e si calcola il punteggio $s_{k,j}$. L'importanza è
\[
i_j=s-\frac{1}{K}\sum_{k=1}^{K}s_{k,j}.
\]
Si calcola con \texttt{permutation\_importance(modello, X, y, n\_repeats=K, scoring=...)}; il risultato contiene \texttt{importances\_mean} e \texttt{importances\_std}. Funziona con qualsiasi modello, non solo con le foreste.
\end{tecnica}

\par\noindent\begin{minipage}{\linewidth}
\fig[\linewidth]{importanze.pdf}
\didascalia{Le 8 feature più importanti secondo la MDI nel modello dell'esempio 2 (Breast Cancer), confrontate con l'importanza per permutazione sul test set (AUC, 20 ripetizioni). Valori reali dall'esecuzione.}
\end{minipage}\par

\begin{esempio}[Cosa ci dice il confronto]
Le due misure concordano sulle feature in cima (\texttt{worst area}, \texttt{worst perimeter}, \texttt{worst concave points}), ma la permutazione dà valori \hl{molto piccoli}: mescolare \texttt{worst area} fa calare l'AUC solo di circa $0.004$. Il motivo è che in questo dataset molte feature sono \key{fortemente correlate} (raggio, perimetro e area misurano quasi la stessa cosa): se ne mescoli una, il modello usa le altre e quasi non se ne accorge.
\end{esempio}

\begin{attenzione}[Feature correlate]
Con feature correlate l'importanza per permutazione \hl{sottostima} ciascuna di esse, e la MDI divide il merito tra loro in modo arbitrario. Nessuna delle due misure dice che una variabile «causa» il target. Una soluzione suggerita dalla documentazione di scikit-learn è raggruppare le feature correlate (ad esempio con un clustering gerarchico sulle correlazioni) e tenerne una per gruppo.
\end{attenzione}

\begin{consiglio}[Quale usare]
Per un'occhiata veloce va bene \texttt{feature\_importances\_}. Per una relazione, una tesi o una decisione usa \texttt{permutation\_importance} su un \key{set di validazione o test}, riportando media e deviazione standard.
\end{consiglio}

\begin{appunti}{Riepilogo del capitolo}
\begin{itemize}
  \item MDI = somma dei cali di impurità dovuti alla feature, media sugli alberi, normalizzata a 1.
  \item MDI: veloce, ma calcolata sul training e distorta verso feature ad alta cardinalità.
  \item Permutazione: $i_j=s-\frac1K\sum_k s_{k,j}$, su dati non visti; più affidabile.
  \item Feature correlate: entrambe le misure vanno interpretate con cautela; importanza $\ne$ causalità.
\end{itemize}
\end{appunti}
